% ============================================================
\chapter{Applications}
\label{ch:applications}
% ============================================================

\begin{center}
\textit{``Theory is when you know everything but nothing works.
Practice is when everything works but nobody knows why.''}\\[4pt]
--- attributed to Albert Einstein
\end{center}

\bigskip

This chapter presents major applications of stochastic calculus:
the Black--Scholes formula in finance, the Feynman--Kac theorem linking
PDEs and SDEs, Kalman filtering, applications in biology, and Monte
Carlo methods for SDEs.

% ------------------------------------------------------------
\section{The Black--Scholes Formula}
\label{sec:black_scholes}
% ------------------------------------------------------------

\begin{definition}[Black--Scholes Model]
\index{Black--Scholes}
In the \textbf{Black--Scholes model}, the price of a risky asset follows
geometric Brownian motion:
\[
  dS_t = \mu\,S_t\,dt + \sigma\,S_t\,dB_t,
\]
and there exists a risk-free asset with return $r$. One assumes
no arbitrage opportunities and continuous frictionless trading.
\end{definition}

\begin{theorem}[Black--Scholes Formula]
\label{thm:bs}
\index{Black--Scholes formula}
The price at time $t$ of a European call option with strike $K$ and
maturity $T$ is:
\[
  C(t, S_t) = S_t\,\Phi(d_1) - K\,e^{-r(T-t)}\,\Phi(d_2),
\]
where $\Phi$ is the standard normal CDF and:
\[
  d_1 = \frac{\ln(S_t/K) + (r + \sigma^2/2)(T-t)}{\sigma\sqrt{T-t}},
  \qquad d_2 = d_1 - \sigma\sqrt{T-t}.
\]
\end{theorem}

\begin{proof}
Under the risk-neutral measure $\mathbb{Q}$, the discounted price is
a martingale: $C(t, S_t) = e^{-r(T-t)}\,\E^{\mathbb{Q}}[(S_T - K)^+
\mid \mathcal{F}_t]$. Under $\mathbb{Q}$,
$S_T = S_t\exp[(r - \sigma^2/2)(T-t) + \sigma(B_T^{\mathbb{Q}} -
B_t^{\mathbb{Q}})]$ where $B^{\mathbb{Q}}$ is a $\mathbb{Q}$-Brownian
motion (Girsanov's theorem\index{Girsanov}). Setting
$Z = (B_T^{\mathbb{Q}} - B_t^{\mathbb{Q}})/\sqrt{T-t} \sim \mathcal{N}(0,1)$:
\[
  C = e^{-r(T-t)}\int_{-\infty}^{\infty}
    \bigl(S_t\,e^{(r-\sigma^2/2)(T-t)+\sigma\sqrt{T-t}\,z} - K\bigr)^+
    \frac{e^{-z^2/2}}{\sqrt{2\pi}}\,dz.
\]
The integral splits over $\{z > -d_2\}$. Substituting $z' = z -
\sigma\sqrt{T-t}$ in the first term yields the stated expressions.
\end{proof}

\begin{intuition}
The Black--Scholes formula shows that the option price does \emph{not}
depend on the expected return $\mu$ of the asset, only on its volatility
$\sigma$. This is a consequence of no-arbitrage: the option can be
perfectly hedged by dynamically adjusting a portfolio of the stock and
the risk-free asset.
\end{intuition}

\begin{definition}[Black--Scholes PDE]
\index{Black--Scholes PDE}
The pricing function $V(t,S)$ satisfies the PDE:
\[
  \frac{\partial V}{\partial t} + r\,S\,\frac{\partial V}{\partial S}
  + \frac{1}{2}\sigma^2 S^2\,\frac{\partial^2 V}{\partial S^2}
  - r\,V = 0,
\]
with terminal condition $V(T,S) = (S-K)^+$ for a call.
\end{definition}

\begin{keyformulas}
\textbf{Black--Scholes Formulas and Greeks}
\begin{align*}
  C &= S\,\Phi(d_1) - Ke^{-r\tau}\Phi(d_2), \quad \tau = T-t \\[4pt]
  \Delta &= \frac{\partial C}{\partial S} = \Phi(d_1), \qquad
  \Gamma = \frac{\partial^2 C}{\partial S^2}
    = \frac{\varphi(d_1)}{S\sigma\sqrt{\tau}} \\[4pt]
  \Theta &= \frac{\partial C}{\partial t}
    = -\frac{S\sigma\varphi(d_1)}{2\sqrt{\tau}}
    - rKe^{-r\tau}\Phi(d_2), \qquad
  \mathcal{V} = \frac{\partial C}{\partial\sigma}
    = S\sqrt{\tau}\,\varphi(d_1)
\end{align*}
\end{keyformulas}

% ------------------------------------------------------------
\section{The Feynman--Kac Theorem}
\label{sec:feynman_kac}
% ------------------------------------------------------------

\begin{theorem}[Feynman--Kac]
\label{thm:feynman_kac}
\index{Feynman--Kac}
Let $u(t,x)$ solve the PDE:
\[
  \frac{\partial u}{\partial t} + b(x)\,\frac{\partial u}{\partial x}
  + \frac{1}{2}\sigma^2(x)\,\frac{\partial^2 u}{\partial x^2}
  - r\,u = 0, \qquad u(T,x) = g(x).
\]
Then $u$ admits the probabilistic representation:
\[
  u(t,x) = \E\bigl[e^{-r(T-t)}\,g(X_T) \mid X_t = x\bigr],
\]
where $X$ solves $dX_s = b(X_s)\,ds + \sigma(X_s)\,dB_s$.
\end{theorem}

\begin{proof}[Proof sketch]
Apply It\^o's formula to $e^{-rs}\,u(s,X_s)$ on $[t,T]$. The $ds$ term
vanishes because $u$ satisfies the PDE, leaving a stochastic integral
(a martingale). Taking the conditional expectation, the stochastic
integral drops out, giving
$e^{-rt}\,u(t,x) = \E[e^{-rT}\,g(X_T) \mid X_t = x]$.
\end{proof}

\begin{remark}
The Feynman--Kac theorem establishes a fundamental link between
parabolic PDEs and diffusion processes. It allows solving PDEs via Monte
Carlo simulation, and conversely, analyzing stochastic problems using
PDE techniques.
\end{remark}

% ------------------------------------------------------------
\section{Kalman Filtering}
\label{sec:kalman}
% ------------------------------------------------------------

\begin{definition}[Linear Filtering Model]
\index{Kalman filter}
Consider the continuous-time linear system:
\begin{align*}
  dX_t &= A\,X_t\,dt + \sigma_X\,dB_t^{(1)} \quad \text{(signal)}, \\
  dY_t &= H\,X_t\,dt + \sigma_Y\,dB_t^{(2)} \quad \text{(observation)},
\end{align*}
where $B^{(1)}$ and $B^{(2)}$ are independent Brownian motions.
The \emph{filtering problem} is to estimate $X_t$ from the observations
$(Y_s)_{0 \leq s \leq t}$.
\end{definition}

\begin{theorem}[Kalman--Bucy Filter]
\label{thm:kalman_bucy}
\index{Kalman--Bucy}
The optimal estimator $\hat{X}_t = \E[X_t \mid \mathcal{F}_t^Y]$
satisfies the SDE:
\[
  d\hat{X}_t = A\,\hat{X}_t\,dt
    + \frac{P_t\,H^T}{\sigma_Y^2}\,\bigl(dY_t - H\,\hat{X}_t\,dt\bigr),
\]
where $P_t = \E[(X_t - \hat{X}_t)^2 \mid \mathcal{F}_t^Y]$ is the
error variance, solving the Riccati equation:
\[
  \frac{dP_t}{dt} = 2A\,P_t + \sigma_X^2 - \frac{P_t^2\,H^2}{\sigma_Y^2}.
\]
\end{theorem}

\begin{intuition}
The Kalman filter achieves an optimal trade-off between the dynamic
model prediction ($A\,\hat{X}_t\,dt$) and the innovation from
observations ($dY_t - H\,\hat{X}_t\,dt$). The Kalman gain
$K_t = P_t H^T/\sigma_Y^2$ weighs these two sources: the larger the
uncertainty $P_t$, the more weight is given to observations.
\end{intuition}

% ------------------------------------------------------------
\section{Applications in Biology}
\label{sec:biology}
% ------------------------------------------------------------

\begin{definition}[Stochastic Gene Expression]
\index{stochastic gene expression}
Gene expression is an intrinsically stochastic process. The protein
count $P_t$ can be modeled by the Langevin-type SDE:
\[
  dP_t = (k_{\text{prod}} - k_{\text{deg}}\,P_t)\,dt
    + \sqrt{k_{\text{prod}} + k_{\text{deg}}\,P_t}\,dB_t,
\]
where $k_{\text{prod}}$ is the production rate and $k_{\text{deg}}$ the
degradation rate.
\end{definition}

\begin{proposition}[Gene Expression Noise]
At equilibrium, the mean protein count is
$\bar{P} = k_{\text{prod}}/k_{\text{deg}}$ and the \emph{coefficient
of variation} (noise measure) is:
\[
  \text{CV}^2 = \frac{\Var(P)}{\E[P]^2} \approx \frac{1}{\bar{P}}.
\]
The noise is thus more significant when the mean molecule count is low.
\end{proposition}

\begin{remark}
This relation $\text{CV}^2 \approx 1/\bar{P}$ is known as
\emph{Poissonian noise}\index{Poissonian noise}. It was experimentally
verified in bacterial cells by Elowitz et al.\ (2002).
\end{remark}

% ------------------------------------------------------------
\section{Monte Carlo Methods for SDEs}
\label{sec:mc_sde}
% ------------------------------------------------------------

\begin{definition}[Euler--Maruyama Scheme]
\index{Euler--Maruyama}
The \textbf{Euler--Maruyama scheme} for the SDE
$dX_t = b(X_t)\,dt + \sigma(X_t)\,dB_t$ is:
\[
  X_{n+1} = X_n + b(X_n)\,\Delta t + \sigma(X_n)\,\sqrt{\Delta t}\,Z_n,
  \qquad Z_n \sim \mathcal{N}(0,1)\;\text{i.i.d.}
\]
\end{definition}

\begin{theorem}[Strong Convergence of Euler--Maruyama]
Under Lipschitz conditions on $b$ and $\sigma$, the Euler--Maruyama
scheme converges with strong order $1/2$:
\[
  \E\left[\sup_{0 \leq t \leq T}
  \abs{X_t - X_t^{(\Delta t)}}^2\right]^{1/2} \leq C\,\sqrt{\Delta t}.
\]
\end{theorem}

\begin{definition}[Milstein Scheme]
\index{Milstein}
The \textbf{Milstein scheme} improves Euler--Maruyama by adding the
It\^o--Taylor correction:
\[
  X_{n+1} = X_n + b(X_n)\,\Delta t + \sigma(X_n)\,\sqrt{\Delta t}\,Z_n
    + \frac{1}{2}\,\sigma(X_n)\,\sigma'(X_n)\,(Z_n^2 - 1)\,\Delta t.
\]
This scheme achieves strong order 1:
$(\E[\sup_t \abs{X_t - X_t^{(\Delta t)}}^2])^{1/2} \leq C\,\Delta t$.
\end{definition}

\begin{attention}
For option pricing (weak convergence), Euler--Maruyama already has weak
order 1, so Milstein offers no improvement for computing $\E[g(X_T)]$.
The distinction between strong and weak convergence is crucial in
practice.
\end{attention}

\begin{keyformulas}
\textbf{Numerical Schemes for SDEs}
\begin{align*}
  \text{Euler--Maruyama:}\quad & X_{n+1} = X_n + b\,\Delta t
    + \sigma\sqrt{\Delta t}\,Z_n
    & \text{(strong order } 1/2\text{)} \\[4pt]
  \text{Milstein:}\quad & X_{n+1} = X_n + b\,\Delta t
    + \sigma\sqrt{\Delta t}\,Z_n
    + \tfrac{1}{2}\sigma\sigma'(Z_n^2-1)\Delta t
    & \text{(strong order } 1\text{)} \\[4pt]
  \text{Monte Carlo:}\quad & \E[g(X_T)] \approx \frac{1}{M}\sum_{m=1}^M
    g(X_T^{(m)})
    & \text{(error } O(1/\sqrt{M})\text{)}
\end{align*}
\end{keyformulas}

\begin{minted}{python}
import numpy as np

def euler_maruyama(b, sigma, x0, T, n_steps, n_paths):
    """Euler-Maruyama scheme for a scalar SDE."""
    dt = T / n_steps
    sqrt_dt = np.sqrt(dt)
    X = np.zeros((n_paths, n_steps + 1))
    X[:, 0] = x0
    for k in range(n_steps):
        Z = np.random.randn(n_paths)
        X[:, k+1] = (X[:, k] + b(X[:, k]) * dt
                     + sigma(X[:, k]) * sqrt_dt * Z)
    return X

# Example: geometric Brownian motion
mu_val, sigma_val, S0 = 0.05, 0.2, 100.0
paths = euler_maruyama(
    b=lambda x: mu_val * x,
    sigma=lambda x: sigma_val * x,
    x0=S0, T=1.0, n_steps=1000, n_paths=10000
)
# Call option price by Monte Carlo
K = 105.0
r = 0.03
call_price = np.exp(-r) * np.mean(np.maximum(paths[:, -1] - K, 0))
\end{minted}

\begin{exercise}
Derive the Black--Scholes formula for a put option using put-call
parity: $C - P = S - Ke^{-r(T-t)}$. Verify that
$P = Ke^{-r\tau}\Phi(-d_2) - S\Phi(-d_1)$.
\end{exercise}

\begin{exercise}
Apply the Feynman--Kac theorem to show that the price of a knock-out
barrier option satisfies the Black--Scholes PDE with boundary condition
$V(t, B) = 0$ on the barrier $S = B$.
\end{exercise}

\begin{exercise}
Implement the discrete Kalman filter to track a moving object observed
with noise. Compare performance with and without a velocity model.
\end{exercise}

\begin{exercise}
Compare the Euler--Maruyama and Milstein schemes for GBM: simulate
$M = 10{,}000$ paths, estimate $\E[S_T]$, and compare the numerical
error for different step sizes $\Delta t$. Verify the theoretical
convergence orders.
\end{exercise}

\begin{exercise}
Simulate a stochastic gene expression model with
$k_{\text{prod}} = 10$, $k_{\text{deg}} = 0.1$ using the Gillespie
algorithm. Verify that the stationary distribution is approximately
Poisson and compute the coefficient of variation.
\end{exercise}
